\scanpage{95}
\mspage{15}
\deleted{plus} simples, qui pour $I'=p^{-1}(I)$ ne dépend pas du choix de $I'$ i.e. de $I$. Si on veut que l'action remontée $\Theta'$ soit l'identité sur $I'$, on a unicité, et la condition supplémentaire que $\widetilde{\Theta}$ doit bien opérer \uline{trivialement} sur $I'$.

Revenons à l'action de $\Gamma$ \inserted{$=\operatorname{Gal}\overline{K}/K$} sur un $\mathcal{T}=\deleted{\unsure{\Pi_1}}\widehat{\Pi}_1(X,I)$ avec $I\subset X_K$, on trouve donc une condition nécessaire pour qu'un \unsure{autom.} $\Theta$ de $\mathcal{T}$ \inserted{trivial sur les objets} provienne d'un $u\in\Gamma$:

\margin{Ces deux conditions ne dépendent en fait que de l'opération de \deleted{$\unsure{\Theta}$} $\Theta$ sur le groupe fondamental en \uuline{un} point $x_0\in I$.}
a) Pour tout rev. \inserted{étale} fini $X'$ de $X$, \inserted{(défini sur $K$)} \deleted{\unsure{il considère l'action de}} \inserted{\unsure{il y a} une action $\widetilde{\Theta}$} de $\Theta$ sur $\widehat{\Pi}_1(X',I')$, où $I'=X'|I$, qui remonte $\Theta$, et qui est l'identité sur $I'\cap X'_K$.

b) De plus, pour $\forall$ hom. (défini sur $K$)
\[
X'\xrightarrow{f}Y
\]
et tout $x'\in I'\cap X'_K$, le noyau de
\[
(*)\qquad\widehat{\pi}_1(X',x')\longrightarrow\widehat{\pi}_1(\unsure{X|Y},f(x')),
\]
est stable sous $\unsure{\widetilde{\Theta}^{\prime}|\widetilde{\Theta}}$ \inserted{(pour $\widetilde{\Theta}$ convenable!)} --- plus généralement, sans supposer $x'\in X'_K$, mais choisissant $\nu\in\unsure{\mathbb{N}|\mathbf{N}}^{+}$ tel que $\Theta^\nu x'=x'$, il faut que \deleted{$\unsure{\Theta}$} le noyau de $(*)$ soit stable par $\Theta^\nu$.

\margin{\uline{NB} Ces conditions \unsure{deviennent informulables}, tant que l'on se \unsure{bornait} au point de vue des actions \uline{extérieures} de $\Gamma$!}
Ce sont là des conditions draconiennes. Je pense que ce n'est pas trop m'avancer,
\mspage{94}
